\subsection{Испарение и конденсация}
%%%%%%%%%%%%%%%%%%%%%%%%%%%%%%%%%%%%%%%%%%%%%%%%%%%%
%%%%%%%%%%%%%%%%%%%%%%%%%%%%%%%%%%%%%%%%%%%%%%%%%%%%
\z{Израсходовав $m_{\text{б}} = 800\,\text{г}$ бензина, воду массой $m_{\text{в}} = 50\,\text{кг}$ нагрели от $t_1 = 20\,\degree\text{C}$ до $t_2 = 100\,\degree\text{C}$ и часть воды выпарили. Какая масса образовавшегося пара, если $\eta = 60\%$ теплоты сгорания бензина передано воде?}
%
{$m = \frac{\eta q m_{\text{б}} - c_{\text{в}} m_{\text{в}} \Delta t}{L} = 1.8\,\text{кг}$}

%%%%%%%%%%%%%%%%%%%%%%%%%%%%%%%%%%%%%%%%%%%%%%%%%%%%
\z{Сколько необходимо сжечь спирта, чтобы расплавить лед массой $m_{\text{л}} = 2\,\text{кг}$, взятый при температуре $t_1 = -5\,\degree\text{C}$, а полученную воду нагреть до кипения и $m_{\text{исп}} = 1\,\text{кг}$ воды превратить в пар? КПД спиртовки $\eta = 40\%$.}
%
{$m = \frac{c_{\text{л}} m_{\text{л}} \Delta t + \lambda m + c_{\text{в}} m_{\text{в}} \Delta t + L m_{\text{п}}}{\eta q} = 0.32\,\text{кг}$}

%%%%%%%%%%%%%%%%%%%%%%%%%%%%%%%%%%%%%%%%%%%%%%%%%%%%
\z{В калориметр, содержащий воду массой $m_{\text{в}} = 500\,\text{г}$ при температуре $t_{\text{в}} = 20\,\degree\text{C}$, впустили водяной пар при температуре $t_{\text{п}} = 100\,\degree\text{C}$. Какая температура $t$ установится в калориметре, если масса пара равна $m_{\text{п1}} = 10\,\text{г}$? А если $m_{\text{п2}} = 100\,\text{г}$? Какой станет масса воды $m$ в каждом из этих случаев?}
%
{\puan $T_{\text{к}} = \frac{L m_{\text{п}} + c_{\text{в}} (m_{\text{п}} t_{\text{п}} + m_{\text{в}} t_{\text{в}})}{c_{\text{в}} (m_{\text{п}} + m_{\text{в}})} = 32\,\degree\text{C},\ m = 510\,\text{г}$; 
\par
\puan $T_{\text{к}} = 100\,\degree\text{C},\ m = \frac{c_{\text{в}} m_{\text{в}} (t_{\text{п}} - t_{\text{в}})}{L} + m = 0.573\,\text{кг}$}

%%%%%%%%%%%%%%%%%%%%%%%%%%%%%%%%%%%%%%%%%%%%%%%%%%%%
\z{До какой температуры нагреется вода объемом $V = 0.8\,\text{л}$, находящаяся в медном калориметре массой $m_{\text{к}} = 700\,\text{г}$ и имеющая температуру $t_1 = 12\,\degree\text{C}$, если впустить в калориметр водяной пар массой $m_{\text{п}} = 50\,\text{г}$ при температуре $t_{\text{п}} = 100\,\degree\text{C}$?}
%
{$t = \frac{t_{\text{в}} (c_{\text{в}} m_{\text{в}} + c_{\text{м}} m_{\text{м}}) + m_{\text{п}} (L + c_{\text{в}} t_{\text{к}})}{c_{\text{в}} m_{\text{в}} + c_{\text{м}} m_{\text{м}} + c_{\text{в}} m_{\text{п}}} = 46.7\,\degree\text{C}$}

%%%%%%%%%%%%%%%%%%%%%%%%%%%%%%%%%%%%%%%%%%%%%%%%%%%%
\z{В кастрюлю налили холодной воды при температуре $t_1 = 10\,\degree\text{C}$ и поставили на электроплитку. Через $\tau_1 = 10\,\text{мин}$ вода закипела. Через какое время после этого она выкипит полностью? Масса кастрюли много меньше массы воды.}
%
{$T = \frac{T q}{c \Delta t} = 3650\,\text{с}$}

%%%%%%%%%%%%%%%%%%%%%%%%%%%%%%%%%%%%%%%%%%%%%%%%%%%%
\z{В калориметр, содержащий $m_{\text{в}} = 1\,\text{кг}$ воды, впустили водяной пар массой $m_{\text{п}} = 40\,\text{г}$, имеющий температуру $t_{\text{п}} = 100\,\degree\text{C}$. Какой была начальная температура воды, если конечная температура в калориметре оказалась равной $t_{\text{к}} = 60\,\degree\text{C}$?}
%
{$t = t_0 - \frac{m_1 \lambda}{c m_2} - \frac{m_1 \Delta t}{m_2} = 36.5\,\degree\text{C}$}

%%%%%%%%%%%%%%%%%%%%%%%%%%%%%%%%%%%%%%%%%%%%%%%%%%%%
\z{В калориметре находится лед массой $m_{\text{л}} = 500\,\text{г}$ при температуре $t_{\text{л}} = -10\,\degree\text{C}$. Какая температура установится в калориметре, если в него впустить пар массой $m_{\text{п}} = 80\,\text{г}$, имеющий температуру $t_{\text{п}} = 100\,\degree\text{C}$.}
%
{$t = \frac{L m_{\text{п}} + 100 c_{\text{в}} m_{\text{п}} - 10 c_{\text{л}} m_{\text{л}} - \lambda m_{\text{л}}}{c_{\text{в}} (m_{\text{л}} + m_{\text{п}})} = 15.2\,\degree\text{C}$}

